\section{Introduction}

\subsection{Problem Statement}

Under finite information and finite resolution constraints, ``observation'' cannot be equated with full access to continuous state, but should be viewed as a protocol jointly determined by geometric parameters, dynamical sampling, and readout kernels. If we require this protocol to satisfy: (i) the ontological layer allows realization of non-periodic structures via high-dimensional periodic objects; (ii) observer differences are rigorously parameterizable; (iii) the readout process generates verifiable invariants; (iv) the steady state is given by computable canonical forms; (v) the emergence of time can be interpreted as information refinement rather than an external parameter, then we need a closed deduction chain from geometry--dynamics--information theory to observables.

\subsection{Main Contributions}

Under strict separation of \textbf{definitions}, \textbf{provable theorems}, \textbf{model class assumptions}, and \textbf{interpretation-layer propositions}, this paper accomplishes the following:
\begin{enumerate}
  \item \textbf{(Definition layer)} Construction of a three-dimensional quasiperiodic point set via six-dimensional lattice cut-and-project scheme, parameterizing observers through internal phase, achieving an auditable representation of ``same ontology—multiple realizations.''
  \item \textbf{(Theorem layer)} Derivation of Sturmian mechanical words from torus rotation of a one-dimensional factor and binary partition readout, providing a complexity minimality invariant.
  \item \textbf{(Definition layer)} Construction of Zeckendorf legal language and stabilization kernel (folding), normalizing arbitrary binary readout windows to the stable state space of ``no adjacent $1$'s.''
  \item \textbf{(Theorem layer + model assumption)} Definition of time fibers as preimages of the stabilization kernel, proof that under ergodic and generating partition conditions the exponential contraction rate of observed history cylinder measures is given by metric entropy, thereby anchoring the ``arrow of time'' to a computable information rate.
  \item \textbf{(Model definition)} Characterization of observation-reachable closure via node--fiber incidence graphs, providing ``mass-like/delay-like'' observable function sub-bundles induced by fiber statistics.
  \item \textbf{(Interpretation-layer proposition)} Proposal of visibility projection and golden ratio fixed point framework, discussion of non-integer effective dimension and arithmetic interfaces, explicitly excluding these from the closed theorem system.
\end{enumerate}

\subsection{Article Structure}

Chapter 2 provides cut-and-project scheme and observer parameterization; Chapters 3--4 establish sampling and Sturmian binary readout under finite resources; Chapter 5 defines Zeckendorf stabilization kernel; Chapter 6 formalizes time as historical refinement of preimage fibers and provides information rate characterization; Chapter 7 provides Levi incidence graph closure and observable function sub-bundles; Chapter 8 discusses statistical equivalence under ergodicity; Chapter 9 provides multiscale and effective dimension discussion; Chapter 10 provides reproducible protocol; appendices provide proof outlines, algorithms, and extension interfaces.
